\RequirePackage{silence}
\WarningsOff % Turn off all warnings from all packages
\ErrorsOff   % (Caution) Use only if you are sure the code has no severe logic errors
\documentclass[a4paper,14pt, twoside]{memoir}
\usepackage[utf8x]{inputenc}
\usepackage[english]{babel}    
\usepackage{url}
\usepackage{microtype}
\usepackage{color}
\usepackage{amsmath}
\usepackage{amsthm}         % Provides proof environment
\usepackage{amssymb}
\usepackage{mathtools}     % Includes amsmath but optimized and with more features
\usepackage{hyperref}
\usepackage{enumitem}
\usepackage{autobreak}
\usepackage{array}          % Support for notation table formatting
\usepackage{booktabs}       % Professional table rules

\hypersetup{
	colorlinks=true,
	linkcolor=blue,
	filecolor=magenta,
	urlcolor=cyan
}

% Theorem Environments
\newtheorem{definition}{Definition}[chapter]
\newtheorem{md}{Proposition}[chapter]
\newtheorem{theorem}{Theorem}[chapter]
\newtheorem{proposition}{Proposition}[chapter]
\newtheorem{corollary}{Corollary}[chapter]
\newtheorem{remark}{Remark}[chapter]
\newtheorem{example}{Example}[chapter]

% Declare Math Operators
\DeclareMathOperator{\Tr}{Tr}
\DeclareMathOperator{\Sym}{Sym}
\DeclareMathOperator{\setu}{s}
\DeclareMathOperator{\dimU}{\dim}
\DeclareMathOperator{\w}{w}
\DeclareMathOperator{\card}{card}
\DeclareMathOperator{\ucard}{ucard}
\DeclareMathOperator{\pr}{pr}
\DeclareMathOperator{\pgr}{pgr}
\DeclareMathOperator{\scale}{scale}
\DeclareMathOperator{\rad}{rad}
\DeclareMathOperator{\Val}{Val}
\DeclareMathOperator{\Conf}{Conf}

% Custom commands
\newcommand{\set}[1]{\left\{#1\right\}} % set
\newcommand{\tbd}[1]{\mathcal{#1}} % uncertain set
\newcommand{\s}[1]{\set{#1}_\mathrm{u}} % uncertain number
\newcommand{\kbd}[1]{{#1}_\mathrm{u}} % uncertain interval
\newcommand{\bd}[1]{\left(#1\right)_\mathrm{u}} % uncertain expression
\newcommand{\tsbd}[1]{\mathbf{U}(\mathbb{#1})}
\newcommand{\tsbdcapn}[2]{\mathbf{U^{#2}}(\mathbb{#1})}
\newcommand{\tsbdts}[1]{\mathbf{U}_{\w}(\mathbb{#1})} % set of uncertain numbers
\newcommand{\ddclmd}{\mathcal{V}} % truth measure of a proposition
\newcommand{\ts}[1]{\w_{#1}} % Weight
\newcommand{\gtcl}[2]{\left(#1\right)_{#2}} % truth measure
\newcommand{\gr}[1]{\text{gr}(#1)} % graph of a function
\newcommand{\hpw}[2]{#1_1\left(#2\right)} % point-wise function
\newcommand{\hbd}[2]{#1_m\left(#2\right)} % uncertain function
\newcommand{\eq}[1]{\left[#1\right]_{eq}} % eqs operator
\newcommand{\kgdd}[1]{\left(#1\right)_1} % pointwise operator
\newcommand{\kgddmr}[1]{\left(#1\right)_{1'}} % real operator
\newcommand{\kgdt}[1]{\left(#1\right)_m} % pointwise operator
\newcommand{\kgdtmr}[1]{\left(#1\right)_{m'}} % pointwise operator

\title{Uncertain Topology}
\author{
	Tran Manh Cuong\thanks{Email: \href{mailto:tmcuong2408@gmail.com}{tmcuong2408@gmail.com} | Phone: 0353237140} \\[0.5em]
	\small Alumnus of Faculty of Mathematics and Informatics (K11)
	\\ Thai Nguyen University of Sciences \\
	\small Address: DJ7 Street, Thoi Hoa, Ho Chi Minh City, Vietnam \\
	\small \textit{Research Field: Mathematics of Uncertainty}
}
\date{Last updated: \today \\ \small DOI / Reference: \url{https://doi.org/10.5281/zenodo.22967037} \cite{uset}}

\begin{document}

\maketitle

\begin{abstract}
	This paper extends Scenario-Based Uncertain Set Theory into the domain of general topology. By avoiding unnecessary symbolic inflation and returning to First Principles, we define an Uncertain Topology simply as a topological structure constructed directly on an Uncertain Universal Space $X$. We show that every open set inherits state-dependent superposition through its Configuration Space $\mathrm{Conf}(A)$. Crucially, we establish the Topological Collapse Theorem: when the uncertain cardinalities of constituent open sets collapse into deterministic integers ($\mathrm{ucard}(A) \in \mathbb{N}$), the uncertain topology degenerates canonical-wise into a classical static topology. From this collapse mechanism, we propose a lightweight, intuitive metric—the Topological Uncertainty Index $\mu_U(\mathcal{T})$—which measures the precise ratio of deterministic capacity over scenario-expanded uncertainty capacity without requiring complex probabilistic densities or fuzzy membership functions.

	\vspace{1em}
	\noindent Keywords: Uncertain Topology, Configuration Space, Topological Collapse, Uncertain Cardinality, Scenario Space, Structural Uncertainty Index, First Principles.
\end{abstract}

\chapter{Uncertain Sets Review}

\begin{definition}[Uncertain Element]
	Let $U$ be a non-empty universal set. An uncertain element $x$ on $U$ is an object whose exact value in $U$ is undetermined, constrained only by a non-empty set of possible values $A \subseteq U$:
	\[
		x \in A
	\]
	The set $A$ is called the possibility domain (or state space) of $x$, denoted by $\mathrm{Val}(x) = A$.
\end{definition}

\begin{definition}[Uncertain Set]
	Let $U$ be a non-empty universal set and $I$ be an index set. An uncertain set $S$ on $U$ is an indexed family of uncertain elements on $U$:
	\[
		S = \{x_i : i \in I\}
	\]
	where each uncertain element $x_i$ ($i \in I$) has a possibility domain $\mathrm{Val}(x_i) \subseteq U$.
\end{definition}

\begin{definition}[Configuration Space]
	Let $S = \{x_i : i \in I\}$ be an uncertain set on $U$. A realized configuration (or empirical state) of $S$ is a classical set $S^* \subseteq U$ obtained when all constituent elements $x_i$ simultaneously realize specific values $v_i \in \mathrm{Val}(x_i)$:
	\[
		S^* = \{v_i : i \in I\}
	\]
	The collection of all classical sets realized by $S$ is called the configuration space of $S$, denoted by $\mathrm{Conf}(S)$.
\end{definition}

\begin{definition}[Uncertain Cardinality]
	Let $S = \{x_i : i \in I\}$ be an uncertain set on $U$. The uncertain cardinality of $S$, denoted by $\mathrm{ucard}(S)$, is an uncertain number on $\mathbb{N}$ whose possibility domain consists of the classical cardinalities of all realized states in its configuration space $\mathrm{Conf}(S)$:
	\[
		\mathrm{Val}(\mathrm{ucard}(S)) = \{ \mathrm{card}(S^*) : S^* \in \mathrm{Conf}(S) \} \subseteq \mathbb{N}
	\]
\end{definition}

\chapter{Uncertain Topology}

\begin{definition}[Uncertain Topology]
	Let $X$ be an uncertain set serving as the universal space. An \emph{Uncertain Topology} on $X$, denoted by $\mathcal{T}$, is a collection of uncertain subsets of $X$ satisfying the standard classical topology axioms across all scenario realizations:
	\begin{enumerate}[label=\roman*)]
		\item $\emptyset \in \mathcal{T}$ and $X \in \mathcal{T}$.
		\item The arbitrary union of open uncertain sets in $\mathcal{T}$ belongs to $\mathcal{T}$.
		\item The finite intersection of open uncertain sets in $\mathcal{T}$ belongs to $\mathcal{T}$.
	\end{enumerate}
	Each element $A \in \mathcal{T}$ is called an \emph{uncertain open set}. Its physical state in any specific scenario is completely characterized by its configuration space $\mathrm{Conf}(A)$.
\end{definition}

\begin{theorem}[Topological Collapse Theorem]
	Let $\mathcal{T}$ be an uncertain topology on an uncertain universe $X$. If for every uncertain open set $A \in \mathcal{T}$, its uncertain cardinality degenerates to a single deterministic natural number:
	\[
		\mathrm{Val}(\mathrm{ucard}(A)) = \{k_A\}, \quad k_A \in \mathbb{N}
	\]
	then the uncertainty across all scenarios vanishes, and $\mathcal{T}$ degenerates canonically into a classical static topology on a constant universe.
\end{theorem}

\begin{proof}
	By definition, $\mathrm{Val}(\mathrm{ucard}(A)) = \{ \mathrm{card}(A^*) : A^* \in \mathrm{Conf}(A) \}$. If $\mathrm{Val}(\mathrm{ucard}(A)) = \{k_A\}$, then every realized configuration $A^* \in \mathrm{Conf}(A)$ has identical classical cardinality $k_A$. This implies that element realization does not induce structural state-collapse or cardinality fluctuation across scenarios. Consequently, $\mathrm{Conf}(A)$ reduces to a singleton configuration set $\{A^*\}$. Thus, every $A \in \mathcal{T}$ behaves identically as a classical set, collapsing $\mathcal{T}$ directly into a classical topology.
\end{proof}

\chapter{Degree of Topological Uncertainty}

\begin{definition}[Topological Uncertainty Index]
	Let $\mathcal{T}$ be an uncertain topology on $X$, and let $\mathrm{Conf}(\mathcal{T})$ be the ensemble of realized classical topologies induced across scenario space. Let $\mathcal{T}_{\mathrm{det}}^*$ be a canonical deterministic baseline configuration of the topology. 

	The \emph{Topological Uncertainty Index} $\mu_U(\mathcal{T}) \in [0, 1)$ is defined as the ratio of deterministic cardinality capacity over total uncertain scenario-expanded cardinality capacity:
	\[
		\mu_U(\mathcal{T}) = 1 - \frac{\sum_{A^* \in \mathcal{T}_{\mathrm{det}}^*} \mathrm{card}(A^*)}{\sum_{A \in \mathcal{T}} \sum_{A^* \in \mathrm{Conf}(A)} \mathrm{card}(A^*)}
	\]
\end{definition}

\begin{proposition}[Properties of $\mu_U(\mathcal{T})$]
	The metric $\mu_U(\mathcal{T})$ satisfies:
	\begin{enumerate}[label=\alph*)]
		\item $\mu_U(\mathcal{T}) = 0$ if and only if $\mathcal{T}$ is a classical static topology.
		\item $\mu_U(\mathcal{T}) \to 1$ as the state-space cardinality fluctuation expands infinitely across scenarios.
	\end{enumerate}
\end{proposition}

\begin{thebibliography}{9}
	\bibitem{uset}
	Tran, M. C. (2026). Scenario-Based Uncertain Set Theory. Zenodo. DOI \href{https://doi.org/10.5281/zenodo.22967037}{https://doi.org/10.5281/zenodo.22967037}
	\bibitem{tran2026extended}
	Tran, M. C. (2026). \textit{Extended Logic and Mathematics of Uncertainty}. Zenodo. DOI: \href{https://doi.org/10.5281/zenodo.22031881}{10.5281/zenodo.22031881}.
	\bibitem{mu}
	tmcuong2408. (2026). tmcuong2408/mu: \textit{Mathematics of Uncertainty Library (Version v0.2.0)} [Computer software]. Zenodo.
	DOI: \href{https://doi.org/10.5281/zenodo.22051784}{10.5281/zenodo.22051784}

\end{thebibliography}

\end{document}